% sec-05: merit review, where the work would run, and the quantitative record.
% Table 20, Figure 12 (diagrams (python)-type clustered infrastructure).
\section{The Merit Review Map}

NSF reviews every proposal against two criteria, Intellectual Merit and Broader
Impacts, and judges SBIR proposals on their commercial impact as well
\cite{nsfpappg,nsf26510}. Table 20 maps each question a reviewer will ask to the
evidence the company can put in front of that reviewer today, and names the
evidence it does not yet have. The last column is as important as the third: a
reviewer trusts a map that shows its own gaps.

\begin{apptable}
\begin{tabularx}{\textwidth}{>{\raggedright\arraybackslash}p{2.5cm} >{\raggedright\arraybackslash}p{3.6cm} >{\raggedright\arraybackslash}p{5.4cm} >{\raggedright\arraybackslash}X}
\headrow \thead{Criterion} & \thead{What NSF asks} & \thead{The evidence today} & \thead{Not yet in hand}\\
\midrule
Intellectual Merit & Advances knowledge? & Model disagreement used as a signal to withhold & Measured rates \\
Intellectual Merit & Well reasoned? & A hard boundary, not instructions; ASME V\&V 40 & An outside review \\
Intellectual Merit & Team qualified? & A dated public record since June 2025 & Staff, on award \\
Intellectual Merit & Resources adequate? & On-premises compute; a NAIRR request & A NAIRR reply \\
Broader Impacts & Benefit to society & Models used only once shown to stay in bounds & Phase I evidence \\
Broader Impacts & Workforce & Eleven planned roles, six needing no license & A first hire \\
Commercial impact & Who pays, and why? & Centers, sponsors and robot makers, for a checkable boundary & A paying customer \\
\end{tabularx}
\end{apptable}
\vspace{-0.60cm}
\tabcap{Table 20. The seven questions an NSF reviewer asks under the three criteria, with the\\
evidence the company can show today and the evidence it does not yet have in hand.}

\subsection{Where compute, data and review would sit}

The resources question is answered by a drawing rather than a sentence. Figure 12
places every component in one of three clusters: the clinical site, where the
model runs and where patient data would stay; the shared NAIRR Pilot resources,
which would see only synthetic cases and the company's code if a request were
granted \cite{nairrpilot2026}; and the public record, which receives the method
and the measured rates. The robot control network is a cluster of its own inside
the site, and the model has no route into it.

\begin{appfloat}[!tb]
\begin{appfig}[diagrams (python)-type / clustered infrastructure, three clusters and one nested]
% Site cluster on the left, with the robot control network nested in its lowest
% row.  NAIRR resources upper right, public record lower right.  Accent tiles run
% on premises; gray tiles are shared or public.  Tile pitch 2.55 cm.
\dgnodew{dgtiled}{0.00}{0.00}{\glyphdb}{Data capture; the model cannot write}{m5}
\dgnodew{dgtiled}{2.55}{0.00}{\glyphai}{Advisory language model, on premises}{m1}
\dgnodew{dgtilem}{5.10}{0.00}{\glyphai}{Two verifier models, other developers}{m2}
\dgnode{dgtile}{0.00}{-2.25}{\glyphlock}{Segmentation: no route to control}{lk}
\dgnode{dgtile}{2.55}{-2.25}{\glyphdoc}{Audit trail, replayable}{m3}
\dgnode{dgtile}{5.10}{-2.25}{\glyphmon}{Surgeon's console, advice as text}{m4}
\dgnodew{dgtilem}{0.00}{-4.55}{\glyphrobot}{Eight-arm platform}{r1}
\dgnodew{dgtilem}{2.55}{-4.55}{\glyphstop}{Stops at 3 ms and 500 ms}{r2}
\dgnodew{dgtilem}{5.10}{-4.55}{\glyphgear}{Force limits, 3 N and 18 N}{r3}
\dgnodeg{dgtileg}{9.90}{0.00}{\glyphcloud}{Allocated compute, if granted}{n1}
\dgnodeg{dgtileg}{12.45}{0.00}{\glyphflask}{Synthetic cases from published values}{n2}
\dgnodeg{dgtileg}{9.90}{-2.25}{\glyphchart}{Measured disagreement rates}{n3}
\dgnodeg{dgtileg}{9.90}{-4.55}{\glyphlink}{Repository and identifiers}{p1}
\dgnodeg{dgtileg}{12.45}{-4.55}{\glyphdoc}{Report of rates and limits}{p2}
\begin{scope}[on background layer]
\node[inner sep=7pt,fit={(r1)(r1l)(r3)(r3l)}] (crx) {};
\node[dgcluster,fit={(m5)(m5l)(m2)(m2l)(crx)},inner sep=9pt] (cs) {};
\node[dgcluster2,fit={(r1)(r1l)(r3)(r3l)}] (cr) {};
\node[dgcluster2,fit={(n1)(n1l)(n2)(n2l)(n3)(n3l)}] (cn) {};
\node[dgcluster2,fit={(p1)(p1l)(p2)(p2l)}] (cp) {};
\end{scope}
\node[dgctitle,anchor=south west] at (cs.north west) {Clinical site, on premises: no patient data leaves it};
\node[dgctitle2,anchor=south east] at ([yshift=1pt]cr.north east) {Isolated robot control network};
\node[dgctitle2,anchor=south west] at (cn.north west) {NAIRR Pilot shared resources};
\node[dgctitle2,anchor=south west] at (cp.north west) {Public record};
\draw[dgedged] (m5) -- node[above,font=\tiny\sffamily,fill=fundwhite,inner sep=1pt] {read} (m1);
\draw[dgedgeb] (m1) -- node[above,font=\tiny\sffamily,fill=fundwhite,inner sep=1pt] {output} (m2);
\draw[dgedgeb] (m2l.south) -- (m4);
\draw[dgedge] (m2l.south west) -- (m3);
\draw[dgedged,-] (m1l.south west) -- (lk);
\draw[dgedged,-] (lkl.south) -- (r1);
\draw[dgedged] (m2) -- node[above,font=\tiny\sffamily,fill=fundwhite,inner sep=1pt] {code only} (n1);
\draw[dgedge] (n2) -- (n1);
\draw[dgedge] (n1l.south) -- (n3);
\draw[dgedge] (n3.east) -| (p2);
\draw[dgedge] (p2) -- (p1);
\legkey{-0.60}{-6.55}{fundaccent}{Runs on premises at the site}
\legkey{6.40}{-6.55}{fundgrayl}{Shared or public; synthetic data only}
\end{appfig}
\vspace{-0.60cm}
\figcaption{Figure 12. Where compute, data and review would sit, with the model and its verifiers\\
on premises, the robot network sealed off, and only code and synthetic cases shared.}
\end{appfloat}

\subsection{The quantitative record behind the robotic Phase 1}

A funder of the verification work is also, one step removed, a funder of the
trial it would make possible, and that trial rests on a dated quantitative
record. In August 2025 the company's ten-arm simulation of the agent returned a
median overall survival of 12.8 months against 5.4 for chemotherapy
\cite{qspmetpancre}. In May 2026 the developer's Phase 3 trial reported 13.2
months against 6.6 in the RAS G12 population of previously treated metastatic
disease \cite{rasolute302}, and the approval that followed reported 13.2 against
6.7 months in the overall population \cite{fdarasonque2026}. The company
describes the closeness of the simulated and reported figures as a chronology,
not a validation, because the populations, regimens and sample sizes differ.
Pancreatic ductal adenocarcinoma remains the cancer with the widest gap between
incidence and survival in the United States \cite{siegel2025statistics}, and
after resection the standard remains combination chemotherapy
\cite{conroy2018folfirinox}.

The program's direct cost is \$700,000 a year for five years, of which \$521,000
is personnel across 3.95 full-time equivalents \cite{movein2026}. The company's
comparable virtual trial work is \textbf{projected} at \$36,330 per run and about
one month, against industry benchmarks above \$120,000 and 4.5 months
\cite{capitalplan2026}. That figure describes simulation work, not the cost of
the clinical trial.
